\documentclass[11pt,a4paper]{article}
\usepackage[margin=0.95in]{geometry}
\usepackage[T1]{fontenc}
\usepackage{lmodern}
\usepackage{amsmath,amssymb}
\usepackage{booktabs}
\usepackage{xcolor}
\usepackage{fancyhdr}
\usepackage{hyperref}
\setlength{\parskip}{0.48em}
\setlength{\parindent}{0pt}
\setlength{\headheight}{15pt}
\linespread{1.08}
\pagestyle{fancy}
\fancyhf{}
\lhead{CST444 Soft Computing - Module 3}
\rhead{Exam-Ready Beginner Guide}
\cfoot{\thepage}
\hypersetup{colorlinks=true,linkcolor=blue,urlcolor=blue}

\newcommand{\examtip}[1]{%
\vspace{0.2em}
\noindent\fcolorbox{black}{gray!12}{\parbox{0.97\linewidth}{\textbf{Exam Tip:} #1}}%
\vspace{0.2em}
}

\begin{document}

\section*{Module 3 Topic Map from OCR Question Papers}
\textbf{Repeated ask areas}
\begin{itemize}
\item Fuzzy set basics, membership function, and fuzzy set operations.
\item Core, support, boundary and properties of fuzzy sets.
\item Plotting membership functions (Age of people, Liquid level).
\item Fuzzy relations, cartesian product, and composition.
\item Max-min and max-product composition numericals.
\item Algebraic sum/product and bounded sum/difference numericals.
\item $\alpha$/$\lambda$ cuts and cut-set properties.
\item Defuzzification methods (max-membership, centroid, weighted average, center of sums).
\end{itemize}

\textbf{High-mark concepts}
\begin{itemize}
\item Defuzzification methods with examples and formulas.
\item Relation construction $R=A\times B$, $P=B\times C$ and $R\circ P$.
\item Cut-set problems and interpretation.
\item Composition methods with full step tables.
\end{itemize}

\textbf{Must-study-first order}
\begin{enumerate}
\item Fuzzy set and membership function meaning.
\item Basic operations and properties.
\item Core-support-boundary and cut sets.
\item Fuzzy relation and composition.
\item Defuzzification methods and method selection.
\item Exam numericals and matrix-style solving.
\end{enumerate}

\section*{Core Theory (Beginner-Friendly, Exam-Formal)}

\subsection*{1) Fuzzy Set and Membership Function}
\textbf{Intuition:} In many real situations, belonging is not only yes/no. Fuzzy logic allows partial belonging.

\textbf{Formal definition}
\[
A=\{(x,\mu_A(x))\mid x\in U\},\qquad \mu_A(x)\in[0,1]
\]
where $U$ is universe of discourse and $\mu_A(x)$ is degree of membership.

\textbf{Discrete representation}
\[
A=\sum_{i=1}^{n}\frac{\mu_A(x_i)}{x_i}
\]

\subsection*{2) Main Features of Membership Function}
\begin{itemize}
\item \textbf{Core:} region where $\mu_A(x)=1$.
\item \textbf{Support:} region where $\mu_A(x)>0$.
\item \textbf{Boundary:} region where $0<\mu_A(x)<1$.
\end{itemize}

\begin{center}
\framebox[0.97\linewidth][c]{
\begin{minipage}{0.93\linewidth}
\vspace{0.08cm}
\centerline{\textbf{Membership-Feature View}}
\centerline{Core: full membership $\rightarrow$ Boundary: partial membership $\rightarrow$ Outside support: zero membership}
\vspace{0.08cm}
\end{minipage}}
\end{center}

\subsection*{3) Fuzzy Set Operations}
For two fuzzy sets $A,B$ on same universe:
\[
\mu_{A\cup B}(x)=\max(\mu_A(x),\mu_B(x)),\quad
\mu_{A\cap B}(x)=\min(\mu_A(x),\mu_B(x)),\quad
\mu_{\bar A}(x)=1-\mu_A(x)
\]

Additional operations frequently asked:
\[
\mu_{A\oplus B}(x)=\mu_A(x)+\mu_B(x)-\mu_A(x)\mu_B(x)\quad\text{(algebraic sum)}
\]
\[
\mu_{A\odot B}(x)=\mu_A(x)\mu_B(x)\quad\text{(algebraic product)}
\]
\[
\mu_{A\boxplus B}(x)=\min(1,\mu_A(x)+\mu_B(x))\quad\text{(bounded sum)}
\]
\[
\mu_{A\boxminus B}(x)=\max(0,\mu_A(x)-\mu_B(x))\quad\text{(bounded difference)}
\]

\subsection*{4) Properties of Fuzzy Sets}
\begin{itemize}
\item Commutative, associative, distributive laws.
\item Idempotent laws: $A\cup A=A$, $A\cap A=A$.
\item Involution: $\overline{\overline{A}}=A$.
\item De-Morgan laws hold in fuzzy form.
\item In fuzzy logic, excluded-middle and contradiction are not strict equalities for all cases.
\end{itemize}

\subsection*{5) Fuzzy Relations and Composition}
A fuzzy relation on $X\times Y$ is
\[
R=\{((x,y),\mu_R(x,y))\mid (x,y)\in X\times Y\}
\]

For $R(X\times Y)$ and $S(Y\times Z)$:
\[
\mu_{R\circ S}(x,z)=\max_{y\in Y}\min\{\mu_R(x,y),\mu_S(y,z)\}\quad\text{(max-min)}
\]
\[
\mu_{R\circ S}(x,z)=\max_{y\in Y}\{\mu_R(x,y)\mu_S(y,z)\}\quad\text{(max-product)}
\]

\subsection*{6) Fuzzification and Defuzzification}
\textbf{Fuzzification:} crisp input $\rightarrow$ fuzzy quantity via membership functions.

\textbf{Defuzzification:} fuzzy output $\rightarrow$ crisp value.

Methods commonly shown in source notes:
\begin{itemize}
\item Max-membership principle.
\item Centroid (center of area).
\item Weighted average method.
\item Center of sums / center of largest area.
\item First of maxima (when multiple maxima exist).
\end{itemize}

\subsection*{7) $\alpha$-Cut and $\lambda$-Cut}
\[
A_\alpha=\{x\in U\mid \mu_A(x)\ge\alpha\}
\]
Strong cut uses $>$; weak cut uses $\ge$.

\examtip{In cut-set numericals, always write the rule first: ``include all elements whose membership is at least the cut level''; then list elements level by level.}

\section*{Numericals and Derivations (Full Steps)}

\subsection*{N1) Intersection, union, complement for given fuzzy sets}
\textbf{Given}
\[
A=\{(x_1,0.5),(x_2,0.1),(x_3,0.9)\},\quad
B=\{(x_1,0.4),(x_2,0.4),(x_3,0.5)\}
\]

\textbf{Find}
$A\cap B$, $A\cup B$, $\bar A$, $\bar B$.

\textbf{Formula}
\[
\mu_{A\cap B}=\min(\mu_A,\mu_B),\quad
\mu_{A\cup B}=\max(\mu_A,\mu_B),\quad
\mu_{\bar A}=1-\mu_A
\]

\textbf{Substitution}
\begin{itemize}
\item $x_1$: min$(0.5,0.4)=0.4$, max$(0.5,0.4)=0.5$.
\item $x_2$: min$(0.1,0.4)=0.1$, max$(0.1,0.4)=0.4$.
\item $x_3$: min$(0.9,0.5)=0.5$, max$(0.9,0.5)=0.9$.
\item Complements: $\bar A=(0.5,0.9,0.1)$, $\bar B=(0.6,0.6,0.5)$.
\end{itemize}

\textbf{Final answer}
\[
A\cap B=\{(x_1,0.4),(x_2,0.1),(x_3,0.5)\}
\]
\[
A\cup B=\{(x_1,0.5),(x_2,0.4),(x_3,0.9)\}
\]

\subsection*{N2) Algebraic sum/product and bounded sum/difference}
\textbf{Given} (module-3 pattern)
\[
A=[1,0.5,0.3,0.2],\quad B=[0.5,0.7,0.2,0.4]
\]

\textbf{Formulas}
\[
A\oplus B=a+b-ab,\ A\odot B=ab,\ A\boxplus B=\min(1,a+b),\ A\boxminus B=\max(0,a-b)
\]

\textbf{Element-wise substitution}
\begin{itemize}
\item Algebraic sum: $[1,0.85,0.44,0.52]$.
\item Algebraic product: $[0.5,0.35,0.06,0.08]$.
\item Bounded sum: $[1,1,0.5,0.6]$.
\item Bounded difference: $[0.5,0,0.1,0]$.
\end{itemize}

\textbf{Sanity check}
All resulting memberships remain in $[0,1]$.

\subsection*{N3) Max-min and max-product composition}
\textbf{Given}
\[
R=\begin{bmatrix}0.6&0.3\\0.2&0.4\end{bmatrix},
\quad
S=\begin{bmatrix}1&0.5&0.3\\0.8&0.4&0.7\end{bmatrix}
\]

\textbf{Max-min formula}
\[
\mu_{T}(x_i,z_k)=\max_j\min\{R_{ij},S_{jk}\}
\]

\textbf{Max-product formula}
\[
\mu_{T}(x_i,z_k)=\max_j\{R_{ij}S_{jk}\}
\]

\textbf{One detailed element each}
\[
T_{11}^{\text{max-min}}=\max(\min(0.6,1),\min(0.3,0.8))=\max(0.6,0.3)=0.6
\]
\[
T_{11}^{\text{max-prod}}=\max(0.6\times1,0.3\times0.8)=\max(0.6,0.24)=0.6
\]

\textbf{Final matrices}
\[
T_{\text{max-min}}=\begin{bmatrix}0.6&0.5&0.3\\0.4&0.4&0.4\end{bmatrix},
\qquad
T_{\text{max-prod}}=\begin{bmatrix}0.6&0.3&0.21\\0.32&0.16&0.28\end{bmatrix}
\]

\subsection*{N4) Defuzzification mini-example (weighted average)}
\textbf{Given}
Rule-output centers $(z_1,z_2,z_3)=(20,40,60)$ and corresponding firing strengths $(0.2,0.5,0.3)$.

\textbf{Formula}
\[
z^*=\frac{\sum_i \mu_i z_i}{\sum_i\mu_i}
\]

\textbf{Substitution}
\[
z^*=\frac{0.2(20)+0.5(40)+0.3(60)}{0.2+0.5+0.3}=\frac{4+20+18}{1}=42
\]

\textbf{Final answer:} defuzzified crisp value is $42$.

\section*{Solved Questions for Module 3 (Self-Contained)}

\subsection*{Part A / 3-Mark Questions (Each Answer Has 6 Scoring Points)}

\textbf{Q1. Plot fuzzy membership function for Age of people (very young, young, middle-aged, old, very old).}
\begin{enumerate}
\item Define universe, for example $U=[0,80]$ years.
\item Choose five overlapping fuzzy labels: very young, young, middle-aged, old, very old.
\item Use triangular/trapezoidal shapes for smooth transition.
\item Ensure neighboring sets overlap to model gradual change.
\item Peak points represent strongest membership for each age label.
\item Final plot must cover whole universe without large gaps.
\end{enumerate}

\textbf{Q2. What are the three basic features of fuzzy membership function?}
\begin{enumerate}
\item Core is region where membership value is exactly 1.
\item Support is region where membership value is greater than 0.
\item Boundary is region where membership is between 0 and 1.
\item Core indicates complete belonging.
\item Support indicates possible belonging.
\item Boundary captures ambiguity/partial belonging.
\end{enumerate}

\textbf{Q3. Define core, support, boundary of a fuzzy set.}
\begin{enumerate}
\item Core$(A)=\{x\mid \mu_A(x)=1\}$.
\item Support$(A)=\{x\mid \mu_A(x)>0\}$.
\item Boundary$(A)=\{x\mid 0<\mu_A(x)<1\}$.
\item Core may be empty for some fuzzy sets.
\item Support can be finite or continuous interval.
\item Boundary region is most important for soft transitions.
\end{enumerate}

\textbf{Q4. List properties of fuzzy set.}
\begin{enumerate}
\item Commutativity for union and intersection.
\item Associativity for union and intersection.
\item Distributive laws hold in fuzzy operations.
\item Idempotent laws: $A\cup A=A$, $A\cap A=A$.
\item Involution: $\bar{\bar A}=A$.
\item De-Morgan laws are valid in fuzzy algebra.
\end{enumerate}

\textbf{Q5. Differentiate fuzzy tolerance and equivalence relation; how convert tolerance to equivalence?}
\begin{enumerate}
\item Both are fuzzy relations on same universe.
\item Tolerance relation is reflexive and symmetric.
\item Equivalence relation is reflexive, symmetric, and transitive.
\item Main difference is transitivity requirement.
\item Conversion is done by transitive closure of tolerance relation.
\item Resulting closure relation becomes fuzzy equivalence relation.
\end{enumerate}

\textbf{Q6. Define composition operation in fuzzy relations.}
\begin{enumerate}
\item Composition combines two compatible fuzzy relations.
\item If $R(X,Y)$ and $S(Y,Z)$, output is relation on $(X,Z)$.
\item Max-min composition uses max of mins.
\item Max-product composition uses max of products.
\item Composition helps build multi-stage fuzzy mappings.
\item It is widely used in fuzzy inference chains.
\end{enumerate}

\textbf{Q7. Given two fuzzy sets, compute algebraic sum/product and bounded sum/difference.}
\begin{enumerate}
\item Use formula $a+b-ab$ for algebraic sum.
\item Use formula $ab$ for algebraic product.
\item Use $\min(1,a+b)$ for bounded sum.
\item Use $\max(0,a-b)$ for bounded difference.
\item Apply element-wise for each universe element.
\item Keep all final values in interval $[0,1]$.
\end{enumerate}

\textbf{Q8. Define fuzzification and defuzzification.}
\begin{enumerate}
\item Fuzzification converts crisp quantity to fuzzy quantity.
\item It assigns membership values to linguistic labels.
\item Defuzzification converts fuzzy output to single crisp value.
\item Max-membership selects point of highest membership.
\item Centroid uses area-balance concept.
\item Weighted average uses weighted centers and firing strengths.
\end{enumerate}

\subsection*{Part B / 14-Mark Questions (Each Answer Has 14 Scoring Points)}

\textbf{Q1. Given fuzzy sets A and B, compute algebraic sum, algebraic product, bounded sum, bounded difference. Then explain defuzzification methods with examples.}
\begin{enumerate}
\item Write all given membership values of $A$ and $B$ in tabular form.
\item State algebraic sum formula $a+b-ab$.
\item Compute algebraic sum element by element.
\item State algebraic product formula $ab$.
\item Compute algebraic product element by element.
\item State bounded sum formula $\min(1,a+b)$.
\item Compute bounded sum element by element.
\item State bounded difference formula $\max(0,a-b)$.
\item Compute bounded difference element by element.
\item Present final four resulting fuzzy sets clearly.
\item Define defuzzification as fuzzy-to-crisp transformation.
\item Explain max-membership principle and when it is used.
\item Explain centroid and weighted-average methods with equations.
\item Conclude with method choice: weighted average for singleton-like outputs, centroid for full-area balance.
\end{enumerate}

\textbf{Q2. Consider fuzzy relation matrix and perform $\lambda$-cut operations for $\lambda=0.9$ and $0^+$. Also compute weighted-average and center-of-sums defuzzification for the given fuzzy output.}
\begin{enumerate}
\item Write relation matrix entries in structured table.
\item For $\lambda$-cut, keep entries with membership $\ge\lambda$.
\item For $\lambda=0.9$, retain only very high-membership cells.
\item For $0^+$ cut, retain all entries with strictly positive membership.
\item Present both cut relations as binary/filtered relation matrices.
\item Define weighted-average formula $z^*=\frac{\sum \mu_i z_i}{\sum\mu_i}$.
\item Identify centers $z_i$ and corresponding membership heights $\mu_i$ from given output set.
\item Substitute all values and compute numerator.
\item Compute denominator and final weighted-average crisp value.
\item Define center-of-sums concept as area-weighted center of clipped outputs.
\item Compute each component area and center location.
\item Apply $z^*=\frac{\sum A_i c_i}{\sum A_i}$.
\item Compare WA and COS outputs briefly.
\item Conclude that different defuzzifiers can produce different but valid crisp outputs.
\end{enumerate}

\textbf{Q3. Plot fuzzy membership functions for liquid level (very small, small, empty, full, very full). Define and explain defuzzification methods.}
\begin{enumerate}
\item Set universe of discourse for liquid level, e.g., 0 to 100 units.
\item Define five linguistic terms requested in the question.
\item Use triangular/trapezoidal curves with overlap.
\item Mark key points (peaks, feet) for each curve.
\item Ensure edge terms cover low and high extremes.
\item Define defuzzification in one formal line.
\item Explain max-membership with formula condition.
\item Explain centroid with integral formula.
\item Explain weighted average with discrete formula.
\item Explain center of sums using area and centroid of each output region.
\item Mention first-of-maxima if multiple peaks exist.
\item State practical preference: centroid is smooth and common.
\item Relate method choice to control precision and computational cost.
\item Conclude with one crisp-value extraction workflow from fuzzy output.
\end{enumerate}

\textbf{Q4. For discrete fuzzy set on $X=\{a,b,c,d,e\}$ find $\alpha$-cut sets for $\alpha=1,0.9,0.6,0^+,0$. Also given A on X and B on Y, find R=A$\times$B, P=B$\times$C and R$\circ$P (max-min).}
\begin{enumerate}
\item Write membership values of the discrete set in table form.
\item Use cut rule $A_\alpha=\{x\mid\mu_A(x)\ge\alpha\}$.
\item Find set at $\alpha=1$ (fully belonging elements).
\item Find set at $\alpha=0.9$.
\item Find set at $\alpha=0.6$.
\item Find weak positive cut $0^+$ (all with membership $>0$).
\item Find $\alpha=0$ cut (entire universe).
\item For fuzzy cartesian product, apply $\mu_R(x,y)=\min(\mu_A(x),\mu_B(y))$.
\item Construct full relation matrix $R=A\times B$.
\item Similarly construct $P=B\times C$.
\item Use max-min composition formula for $R\circ P$.
\item Compute each output cell with ``max of mins'' operation.
\item Present final composed matrix with dimensions $|X|\times|V|$.
\item Conclude interpretation: composition propagates fuzzy compatibility across chained universes.
\end{enumerate}

\textbf{Q5. Define composition in fuzzy relations. Given R and S, compute T using max-min and max-product composition.}
\begin{enumerate}
\item Write given relation matrices $R$ and $S$.
\item Verify compatibility of inner dimension (Y set must match).
\item State max-min formula explicitly.
\item Compute one cell of max-min composition step by step.
\item Compute remaining max-min cells and form matrix.
\item State max-product formula explicitly.
\item Compute one cell of max-product step by step.
\item Compute remaining max-product cells and form matrix.
\item Compare values from both methods in one line.
\item Note: max-product tends to reduce values more aggressively when memberships are small.
\item Verify all output memberships lie in $[0,1]$.
\item Mention computational flow as matrix-style repeated operations.
\item Provide final two matrices clearly labeled.
\item Conclude: choice of composition method affects relation strength in inference pipelines.
\end{enumerate}

\section*{Formula Sheet (Module 3)}
\begin{itemize}
\item Fuzzy set: $A=\{(x,\mu_A(x))\mid x\in U\}$
\item Union: $\mu_{A\cup B}=\max(\mu_A,\mu_B)$
\item Intersection: $\mu_{A\cap B}=\min(\mu_A,\mu_B)$
\item Complement: $\mu_{\bar A}=1-\mu_A$
\item Algebraic sum: $a+b-ab$
\item Algebraic product: $ab$
\item Bounded sum: $\min(1,a+b)$
\item Bounded difference: $\max(0,a-b)$
\item Cartesian product relation: $\mu_{A\times B}(x,y)=\min(\mu_A(x),\mu_B(y))$
\item Max-min composition: $\mu_{R\circ S}(x,z)=\max_y\min(\mu_R(x,y),\mu_S(y,z))$
\item Max-product composition: $\mu_{R\circ S}(x,z)=\max_y(\mu_R(x,y)\mu_S(y,z))$
\item $\alpha$-cut: $A_\alpha=\{x\mid\mu_A(x)\ge\alpha\}$
\item Weighted average defuzzifier: $z^*=\frac{\sum\mu_i z_i}{\sum\mu_i}$
\item Centroid defuzzifier: $z^*=\frac{\int z\mu(z)\,dz}{\int \mu(z)\,dz}$
\end{itemize}

\section*{One-Page Quick Revision Summary}
\textbf{Memory hooks}
\begin{itemize}
\item Fuzzy set = graded membership, not strict yes/no.
\item Core/support/boundary = full/possible/partial regions.
\item Composition = connect two fuzzy relations through common universe.
\item Cut sets = threshold snapshots of fuzzy information.
\item Defuzzification = final crisp decision from fuzzy output.
\end{itemize}

\textbf{Fast exam writing sequence}
\begin{enumerate}
\item Definition in one line.
\item Formula with symbols.
\item Step-by-step substitution.
\item Final set/matrix/value clearly boxed or isolated.
\item One-line interpretation.
\end{enumerate}

\textbf{Last-minute 45-second recall}
\begin{itemize}
\item Use min for intersection and max for union.
\item For $A\times B$, use min of two memberships.
\item Max-min and max-product are not the same.
\item $\alpha$-cut uses threshold rule $\mu\ge\alpha$.
\item Centroid is area-balance method; WA is weighted-center method.
\end{itemize}

\examtip{In matrix composition, show one full cell derivation clearly. Examiners often award method marks quickly when that line is correct.}

\end{document}
